\clearpage
\begin{examplebox}
\textbf{\textcolor{CExample}{例1：求斜率积（基础选择/填空）}}

\textbf{题目：}双曲线\(\frac{x^2}{4}-\frac{y^2}{5}=1\)，
\(A_1(-2,0), A_2(2,0)\)，\(P\)为曲线上非顶点，求\(k_{PA_1}k_{PA_2}\)。

\textbf{速解：}\(a^2=4,b^2=5\)，于是
\[
 \boxed{k_{PA_1}k_{PA_2}=\frac54}.
\]
\textbf{识别信号：}两个端点\(+\)动点\(+\)斜率积，直接用定值。

\medskip
\textbf{\textcolor{CExample}{例2：已知斜率积求离心率}}

\textbf{题目：}双曲线\(\frac{x^2}{a^2}-\frac{y^2}{b^2}=1\)的
两个实轴端点为\(A_1,A_2\)，对曲线上非顶点\(P\)，
已知\(k_{PA_1}k_{PA_2}=3\)，求离心率\(e\)。

\textbf{速解：}
\[
 e^2-1=3\quad\Longrightarrow\quad
 \boxed{e=2}\quad(e>1).
\]
\textbf{易错提醒：}离心率是正数，不能取\(-2\)。

\medskip
\textbf{\textcolor{CExample}{例3：已知斜率和与积，速求点的坐标（重点）}}

\textbf{题目：}在\(C:\frac{x^2}{4}-\frac{y^2}{5}=1\)上，
取\(P\ne A_1,A_2\)，若
\(k_{PA_1}+k_{PA_2}=3\)，求所有点\(P\)的坐标。

\begin{center}
\begin{tikzpicture}[x=1cm,y=1cm,>=stealth,font=\small]
\path[use as bounding box] (-5.4,-3.8) rectangle (5.4,4.1);
\node[font=\bfseries] at (0,3.8) {中心对称：两个斜率交换，乘积不变};
\draw[->,black!50] (-4.6,0)--(4.6,0) node[right] {$x$};
\draw[->,black!50] (0,-3.1)--(0,3.2) node[above] {$y$};
\foreach \s in {-1,1} {\draw[blue!65!black,thick,domain=-3:3,samples=80,variable=\t] plot ({\s*sqrt(4+.8*\t*\t)},\t);}
\draw[orange!85!black,thick] (-2,0)--(3,2.5) (2,0)--(-3,-2.5);
\draw[teal!80!black,thick] (2,0)--(3,2.5) (-2,0)--(-3,-2.5);
\fill (-2,0) circle (2pt) node[above left] {$A_1$};
\fill (2,0) circle (2pt) node[below right] {$A_2$};
\fill (3,2.5) circle (2pt) node[right=8pt] {$P_1$};
\fill (-3,-2.5) circle (2pt) node[left=8pt] {$P_2$};
\node[orange!85!black,fill=BGExplanation] at (-.9,1.35) {$1/2$};
\node[teal!80!black,fill=BGExplanation] at (3.6,1.25) {$5/2$};
\node[teal!80!black,fill=BGExplanation] at (-3.6,-1.25) {$5/2$};
\node[orange!85!black,fill=BGExplanation] at (.9,-1.35) {$1/2$};
\node at (0,-3.6) {$P_1=(3,5/2),\quad P_2=(-3,-5/2),\quad k_1k_2=5/4$};
\end{tikzpicture}

\end{center}
图中两点分别对应两种斜率顺序；因此交换两根的对应不能省略。

\textbf{第一步（把斜率当二次方程的两根）：}
令\(u=k_{PA_1},v=k_{PA_2}\)，则
\(u+v=3,\ uv=\frac54\)，于是
\[
 t^2-3t+\frac54=0
 \quad\Longrightarrow\quad \{u,v\}=\left\{\frac12,\frac52\right\}.
\]
\textbf{第二步（两条直线直接求交，避免繁琐代数）：}
由\(u(x+2)=v(x-2)=y\)，得到通用公式
\[
 \boxed{x=\frac{2(u+v)}{v-u},\qquad
 y=\frac{4uv}{v-u}}\quad(u\ne v).
\]
当\((u,v)=(\frac12,\frac52)\)时，
\((x,y)=(3,\frac52)\)；
当\((u,v)=(\frac52,\frac12)\)时，
\((x,y)=(-3,-\frac52)\)。
\[
 \boxed{P=(3,\tfrac52)\ \text{或}\ P=(-3,-\tfrac52)}.
\]
\textbf{验证：}两个点均满足\(x^2/4-y^2/5=1\)，且相应的斜率和为\(3\)。
\textbf{陷阱：}由韦达解出的两根\textbf{没有自动指定}\(u,v\)对应哪条直线，
需交换对应检查；只写一个点会漏解。

\medskip
\textbf{\textcolor{CExample}{例4：由斜率积反求动点轨迹（逆向第三定义）}}

\textbf{题目：}固定\(A_1(-3,0),A_2(3,0)\)。动点\(P\)到两端点连线的
斜率均存在，且乘积为\(\frac49\)。求轨迹。

\textbf{解：}设\(P(x,y)\)并注意\(x\ne\pm3\)，
\[
 \frac{y^2}{x^2-9}=\frac49
 \Longrightarrow y^2=\frac49(x^2-9)
 \Longrightarrow\boxed{\frac{x^2}{9}-\frac{y^2}{4}=1}.
\]
\textbf{答案：}上式双曲线\textbf{去掉两顶点}\((\pm3,0)\)。
\textbf{原因：}两顶点虽然满足轨迹方程，但斜率条件在顶点处无意义。

\medskip
\textbf{\textcolor{CExample}{例5：纵轴型辨析（高价值变式）}}

\textbf{题目：}已知\(\frac{y^2}{9}-\frac{x^2}{16}=1\)，
两个实轴端点\(B_1(0,-3),B_2(0,3)\)。
曲线上非顶点\(P\)，求\(k_{PB_1}k_{PB_2}\)，并求离心率\(e\)。

\textbf{解：}\(a^2=9,b^2=16\)，纵轴型斜率积为
\[
 \boxed{k_{PB_1}k_{PB_2}=\frac{a^2}{b^2}=\frac9{16}}.
\]
又\(c^2=9+16=25\)，\(\boxed{e=\frac53}\)，
从而\(1/(e^2-1)=9/16\)，与斜率积吻合。
\textbf{易错提醒：}把\(16/9\)直接当成斜率积，属于\textbf{横纵型混用}。

\medskip
\textbf{\textcolor{CExample}{例6：判断“只有斜率和”是否可能（提升变式）}}

\textbf{题目：}仍在\(\frac{x^2}{4}-\frac{y^2}{5}=1\)上，
能否找到非顶点\(P\)，使\(k_{PA_1}+k_{PA_2}=\sqrt5\)？

\textbf{解：}若可行，则\(u+v=\sqrt5\)且\(uv=\frac54\)。
\(u,v\)应为\(t^2-\sqrt5t+\frac54=0\)的实根，判别式
\[
 \Delta=(\sqrt5)^2-4\cdot\frac54=0
 \ \Longrightarrow\ u=v=\frac{\sqrt5}{2}.
\]
然而\(P\ne A_1,A_2\)时，两条顶点连线是\textbf{过两个不同顶点的不同直线}，
其斜率不能相同（否则平行且同过\(P\)，矛盾）。
故\(\boxed{\text{不存在这样的 }P}\)。

\textbf{迁移条件：}若斜率积为\(m>0\)而斜率和为\(s\)，
曲线上非顶点\(P\)存在的充要条件是\(\boxed{s^2>4m}\)，
不是\(s^2\ge4m\)。

\medskip
\textbf{\textcolor{CExample}{例7：从直径所对圆周角认识斜率积 $-1$（初中衔接）}}

\textbf{题目：}固定 $A_1(-2,0),A_2(2,0)$，动点 $P$ 满足两条连线斜率均存在，且
$k_{PA_1}k_{PA_2}=-1$，求轨迹，并给出几何解释。

\textbf{两步速解：}
\[
 \frac{y^2}{x^2-4}=-1\iff\boxed{x^2+y^2=4},
 \qquad P\ne(\pm2,0).
\]
\textbf{几何解：}$k_1k_2=-1\iff PA_1\perp PA_2$；
由圆周角定理，轨迹就是\textbf{以 $A_1A_2$ 为直径的圆，去掉直径端点}。
\textbf{迁移：}这是 C065 的一条\textbf{反向知识入口}，不需要先知道双曲线。

\medskip
\textbf{\textcolor{CExample}{例8：同一对端点，一式统一三种圆锥曲线}}

\textbf{题目：}固定 $A_1(-2,0),A_2(2,0)$；斜率积为常数 $m$。
分别取 $m=-1$、$m=-\frac14$、$m=\frac54$，求轨迹。

\textbf{速解：}总公式 $y^2=m(x^2-4)$，所以
\[
\begin{array}{ll}
 m=-1:& \boxed{x^2+y^2=4}\quad(\text{圆}),\\[3pt]
 m=-\frac14:&\boxed{\dfrac{x^2}{4}+y^2=1}\quad(\text{椭圆}),\\[3pt]
 m=\frac54:&\boxed{\dfrac{x^2}{4}-\dfrac{y^2}{5}=1}\quad(\text{双曲线}).
\end{array}
\]
三条轨迹都\textbf{排除} $A_1,A_2$，否则分母为零、原斜率条件无定义。
\textbf{识别信号：}若题目给两个定点和斜率\textbf{乘积}，先写总公式，再按 $m$ 分类。

\medskip
\textbf{\textcolor{CExample}{例9：用辅助圆核对斜率积（圆幂法的数值演示）}}

\textbf{题目：}$C:\frac{x^2}{4}-\frac{y^2}{5}=1$，$P(3,\frac52)$；
以 $A_1(-2,0),A_2(2,0)$ 为直径作圆，$H(3,0)$ 引切线 $HT$。
用 $PH,HT$ 求斜率积。

\textbf{解：}由切割线定理
\[
 HT^2=HA_1\cdot HA_2=5\times1=5,\qquad PH=\frac52.
\]
因此
\[
 \boxed{k_1k_2=\left(\frac{PH}{HT}\right)^2
 =\frac{25/4}{5}=\frac54}.
\]
\textbf{几何意义：}两个水平底的乘积被\textbf{一个切线长度平方}代替。

\medskip
\textbf{\textcolor{CExample}{例10：半角参数法为什么是“倒数相消”}}

\textbf{题目：}$C:\frac{x^2}{4}-\frac{y^2}{4}=1$；
取 $\theta=\pi/3$，令 $P=(2\sec\theta,2\tan\theta)$。

\textbf{解：}$P=(4,2\sqrt3)$，由半角公式
\[
 k_1=\tan\frac\pi6=\frac1{\sqrt3},\qquad
 k_2=\cot\frac\pi6=\sqrt3,
 \quad\boxed{k_1k_2=1}.
\]
验证 $b^2/a^2=4/4=1$。
\textbf{提醒：}当 $\theta=0$ 时 $P$ 是顶点，斜率条件不成立。

\end{examplebox}
